% ─────────────────────────────────────────────────────────────────────────
% Gaya PANDUAN_ODIN.pdf — dipakai pandoc via --include-in-header
% ─────────────────────────────────────────────────────────────────────────
\usepackage{fontspec}
\usepackage{xcolor}
\usepackage{fancyhdr}
\usepackage{titlesec}
\usepackage{tcolorbox}
\usepackage{enumitem}
\usepackage{newunicodechar}
\usepackage{booktabs}
\usepackage{array}
\usepackage{ragged2e}
\tcbuselibrary{skins,breakable}

% ── Palet ────────────────────────────────────────────────────────────────
\definecolor{odinblue}{HTML}{1B3A5C}
\definecolor{odincyan}{HTML}{00758F}
\definecolor{odingray}{HTML}{4A5568}
\definecolor{odinlight}{HTML}{EDF2F7}
\definecolor{odinred}{HTML}{B03030}
\definecolor{odinamber}{HTML}{9A6700}
\definecolor{odingreen}{HTML}{2F6B4F}
\definecolor{codebg}{HTML}{F7F8FA}
\definecolor{coderule}{HTML}{D8DEE7}

% ── Font: Charter (teks) + Avenir Next (judul) + Menlo (kode) ────────────
\setmainfont{Charter}[Numbers=OldStyle, Ligatures=TeX]
\newfontfamily\headingfont{Avenir Next}[Ligatures=TeX]
\setmonofont{Menlo}[Scale=0.82]
% Menlo memuat panah, box-drawing, dan simbol yang tak ada di Charter.
\newfontfamily\symfont{Menlo}[Scale=0.9]
\newunicodechar{→}{{\symfont→}}
\newunicodechar{←}{{\symfont←}}
\newunicodechar{↓}{{\symfont↓}}
\newunicodechar{▶}{{\symfont▶}}
\newunicodechar{✓}{{\symfont✓}}
\newunicodechar{⚠}{{\symfont⚠}}
\newunicodechar{⚡}{{\symfont⚡}}
\newunicodechar{─}{{\symfont─}}
\newunicodechar{━}{{\symfont━}}
\newunicodechar{│}{{\symfont│}}
\newunicodechar{├}{{\symfont├}}
\newunicodechar{└}{{\symfont└}}
\newunicodechar{┌}{{\symfont┌}}
\newunicodechar{┐}{{\symfont┐}}
\newunicodechar{┘}{{\symfont┘}}
\newunicodechar{≥}{{\symfont≥}}
\newunicodechar{≤}{{\symfont≤}}
\newunicodechar{×}{{\symfont×}}
\newunicodechar{≡}{{\symfont≡}}

% ── Judul bab & seksi ────────────────────────────────────────────────────
\titleformat{\chapter}[display]
  {\headingfont\color{odinblue}}
  {\normalsize\bfseries\color{odincyan}\MakeUppercase{Bab \thechapter}}
  {6pt}{\Huge\bfseries}
  [\vspace{2pt}{\color{odincyan}\titlerule[1.2pt]}]
\titlespacing*{\chapter}{0pt}{-8pt}{22pt}
\titleformat{\section}{\headingfont\Large\bfseries\color{odinblue}}{\thesection}{0.6em}{}
\titleformat{\subsection}{\headingfont\large\bfseries\color{odingray}}{\thesubsection}{0.6em}{}
\titleformat{\subsubsection}{\headingfont\normalsize\bfseries\color{odingray}}{}{0em}{}

% ── Header & footer ──────────────────────────────────────────────────────
\pagestyle{fancy}
\fancyhf{}
\renewcommand{\headrulewidth}{0.4pt}
\renewcommand{\footrulewidth}{0pt}
\fancyhead[L]{\headingfont\footnotesize\color{odingray}Panduan ODIN v2.3.0}
\fancyhead[R]{\headingfont\footnotesize\color{odingray}\nouppercase{\leftmark}}
\fancyfoot[C]{\headingfont\footnotesize\color{odingray}\thepage}
\fancypagestyle{plain}{\fancyhf{}\renewcommand{\headrulewidth}{0pt}%
  \fancyfoot[C]{\headingfont\footnotesize\color{odingray}\thepage}}
\renewcommand{\chaptermark}[1]{\markboth{#1}{}}

% ── Blok kode ────────────────────────────────────────────────────────────
\usepackage{fvextra}
\fvset{breaklines=true, breakanywhere=true, fontsize=\small, commandchars=\\\{\}}
\renewenvironment{Shaded}
  {\begin{tcolorbox}[breakable, colback=codebg, colframe=coderule, boxrule=0.5pt,
     arc=2pt, left=6pt, right=6pt, top=5pt, bottom=5pt, enhanced]}
  {\end{tcolorbox}}

% Blok kode tanpa bahasa dikeluarkan pandoc sebagai `verbatim`, bukan `Shaded`,
% jadi perlu dikotaki terpisah agar tampilannya seragam.
\tcolorboxenvironment{verbatim}{breakable, colback=codebg, colframe=coderule,
  boxrule=0.5pt, arc=2pt, left=6pt, right=6pt, top=5pt, bottom=5pt, enhanced,
  before skip=8pt, after skip=8pt}

% ── Sampul ───────────────────────────────────────────────────────────────
\renewcommand{\maketitle}{%
  \begin{titlepage}
  \thispagestyle{empty}
  \raggedright
  \vspace*{2.2cm}
  {\color{odincyan}\rule{\textwidth}{3pt}}\par
  \vspace{1.1cm}
  {\headingfont\fontsize{76}{80}\selectfont\bfseries\color{odinblue}ODIN}\par
  \vspace{0.7cm}
  {\headingfont\LARGE\color{odinblue}Panduan Lengkap}\par
  \vspace{0.25cm}
  {\headingfont\large\color{odingray}MCP Agent AI untuk Server Linux}\par
  \vspace{1.1cm}
  {\color{odincyan}\rule{4.5cm}{1.2pt}}\par
  \vspace{1.1cm}
  \begin{minipage}{0.88\textwidth}
    \large\itshape\color{odingray}
    Claude Code adalah otak. ODIN adalah tangan.\par\vspace{0.35em}
    Perintah bahasa manusia di laptop, dieksekusi di server Linux yang hidup —
    dengan kartu risiko sebelum setiap perubahan.
  \end{minipage}
  \vfill
  {\color{coderule}\rule{\textwidth}{0.6pt}}\par
  \vspace{0.6cm}
  {\headingfont\large\bfseries\color{odinblue}Versi 2.3.0}\hfill
  {\headingfont\normalsize\color{odingray}September 2026}\par
  \vspace{0.5cm}
  {\headingfont\normalsize\color{odingray}Syamsuddin Ideris\quad·\quad
   syamsuddin.ideris@gmail.com}\par
  \vspace{0.15cm}
  {\headingfont\footnotesize\color{odingray}
   Lisensi MIT\quad·\quad github.com/Syamsuddin/ODIN}\par
  \vspace{0.8cm}
  \end{titlepage}
}

% ── Kutipan → kotak catatan ──────────────────────────────────────────────
\renewenvironment{quote}
  {\begin{tcolorbox}[breakable, colback=odinlight, colframe=odincyan, boxrule=0pt,
     leftrule=3pt, arc=0pt, left=8pt, right=8pt, top=6pt, bottom=6pt, enhanced]}
  {\end{tcolorbox}}

% ── Tabel: lebih lapang, teks rata kiri di kolom lebar ───────────────────
\renewcommand{\arraystretch}{1.25}
\setlength{\tabcolsep}{5pt}

\setlist[itemize]{leftmargin=1.2em, itemsep=2pt, topsep=3pt}
\setlist[enumerate]{leftmargin=1.4em, itemsep=2pt, topsep=3pt}

\usepackage{microtype}
\hyphenpenalty=1200
\emergencystretch=2em
